% ── CAPÍTULO 1: RESUMEN EJECUTIVO ───────────────────────────
\section{Resumen Ejecutivo}
\label{sec:resumen}

\subsection{Descripción del Problema}

El presente proyecto aborda la extensión del problema de la \textbf{Partición de Mínima Información} (MIP, \textit{Minimum Information Partition}) al caso general de $k$-particiones dentro del marco de la \textbf{Teoría de la Información Integrada} (IIT). El problema consiste en determinar cómo dividir un sistema de $n$ variables binarias en $k$ partes independientes de manera que se minimice la pérdida de información integrada. Formalmente, se busca la partición del sistema $V$ en $k$ subconjuntos disjuntos $S_1, S_2, \ldots, S_k$ que minimice la discrepancia entre la dinámica del sistema completo y la dinámica reconstruida desde las partes.

\subsection{Enfoque Algorítmico}

Se implementaron dos estrategias que extienden los algoritmos existentes de bi-particiones:

\begin{itemize}[itemsep=4pt]
    \item \textbf{KGeoMIP}: Extensión del algoritmo geométrico GeoMIP. Aprovecha la representación del espacio de estados como hipercubo $n$-dimensional y una tabla de costos de transiciones para generar y evaluar particiones candidatas en $k$ grupos ($k \in \{2,3,4,5\}$).

    \item \textbf{KQNodes}: Extensión del algoritmo QNodes basado en el método de Queyranne para optimización submodular. Utiliza un oráculo Zeta precomputado y búsqueda exhaustiva o heurística greedy sobre las $\Stirling{n}{k}$ posibles $k$-particiones.
\end{itemize}

Ambas estrategias adaptan automáticamente su modo de búsqueda (exacto vs.\ heurístico) según el número de Stirling de segundo tipo $\Stirling{n}{k}$ respecto a un presupuesto configurable (por defecto $10{,}000$ evaluaciones).

\subsection{Principales Resultados}

Se ejecutaron benchmarks sobre sistemas de 5 a 22 variables (redes N5B, N10A, N10B, N15A, N20A, N21A, N22A). Los experimentos validan:

\begin{itemize}[itemsep=4pt]
    \item Para $k = 2$, ambas extensiones reproducen exactamente los resultados de GeoMIP y QNodes originales (validaciones TV-02 y TV-03).
    \item KQNodes garantiza optimalidad exacta cuando $\Stirling{n}{k} \leq 10{,}000$.
    \item KGeoMIP ofrece tiempos menores con búsqueda heurística guiada geométricamente.
\end{itemize}

\subsection{Limitaciones y Recomendaciones}

La complejidad combinatoria de $\Stirling{n}{k}$ limita la búsqueda exacta a sistemas de tamaño moderado. Para $n \geq 15$ y $k \geq 3$, ambos algoritmos operan en modo heurístico. Se recomienda KQNodes cuando se requiere exactitud en sistemas pequeños, y KGeoMIP cuando prima la velocidad en sistemas grandes.
